\documentclass[11pt]{article}
\usepackage[margin=1in]{geometry}
\usepackage[T1]{fontenc}
\usepackage[utf8]{inputenc}
\usepackage{lmodern}
\usepackage{microtype}
\usepackage{amsmath,amssymb,amsfonts}
\usepackage{booktabs}
\usepackage{siunitx}
\usepackage{xcolor}
\usepackage{setspace}
\usepackage{tcolorbox}
\tcbuselibrary{skins,breakable}

\usepackage{hyperref}
\hypersetup{
    colorlinks=true,
    linkcolor=blue!80!black,
    citecolor=blue!80!black,
    urlcolor=blue!80!black
}

\definecolor{revgray}{rgb}{0.96, 0.96, 0.98}
\definecolor{revborder}{rgb}{0.75, 0.78, 0.85}

\newtcolorbox{reviewercomment}[1][]{
    enhanced,
    breakable,
    colback=revgray,
    colframe=revborder,
    coltitle=black,
    arc=2mm,
    boxrule=0.8pt,
    left=12pt,
    right=12pt,
    top=9pt,
    bottom=9pt,
    fontupper=\itshape,
    before skip=10pt,
    after skip=10pt,
    #1
}

\setlength{\parindent}{0pt}
\setlength{\parskip}{7pt}
\setstretch{1.12}

\begin{document}

\noindent{\LARGE \textbf{Response to Reviewers and Editor}}\\[6pt]
\textbf{Manuscript Number:} SON-D-26-00511\\
\textbf{Date:} \today\\[12pt]

\section*{Overview and Letter to the Editor}

Dear Dr.~Brandes,

Thank you very much for the opportunity to revise and resubmit my manuscript, \textit{Cultural Matching and the Persistence of Social Ties} (SON-D-26-00511), for publication in \textit{Social Networks}. I am deeply grateful to you and the two anonymous reviewers for the thorough, incisive, and constructive guidance provided throughout the review process. The reviewers' comments have pushed me to substantively sharpen my theoretical framework, formalize my empirical methodology, introduce basic structural embeddedness and demographic controls, and perform required sensitivity analyses.

In response to this feedback, I have undertaken a comprehensive revision of the manuscript and my analytical strategy. Below, I provide a point-by-point response to every issue raised by the reviewers. In synthesis, my primary revisions encompass ten key empirical, methodological, and theoretical changes:
(1)~\textit{incorporating dyadic structural embeddedness (triadic closure and common alter counts) into all discrete-time survival models}, confirming that cultural matching protects ties from decay independently of collective triadic clustering;
(2)~\textit{formalizing the discrete-time event history framework and examining the combinatorial distribution of tie sequences across eight panel waves}, including an absorbing first-decay sensitivity model that corroborates my findings;
(3)~\textit{empirically evaluating and documenting the longitudinal stability and test-retest consistency of cultural tastes}, showing that cultural orientations function as durable personal dispositions rather than volatile fads during the undergraduate years;
(4)~\textit{accounting for alters who participate as study egos through cross-classified multilevel models and undirected dyadic clustering}, verifying that two-way clustering does not distort my standard errors or parameter estimates;
(5)~\textit{decomposing cultural matching into active positive endorsement versus shared disinterest}, revealing that mutual positive tastes protect ties from decay;
(6)~\textit{conducting extensive node persistence, panel retention, and within-ego fixed-effects analyses}, ruling out survivorship or panel attrition bias;
(7)~\textit{embedding comprehensive descriptive statistics tables directly within the main text};
(8)~\textit{reframing the theoretical argument away from a confrontational culture-versus-structure dichotomy toward an integrative co-evolutionary and complementary mechanism framework};
(9)~\textit{calibrating my theoretical scope to the developmental dynamics of emerging adulthood and college transitions}; and
(10)~\textit{deepening my theoretical and empirical engagement with campus positional differences, exogenous identity criteria, status asymmetries, and the bounded conversion of cultural capital across social locations}.

Below, I detail each reviewer comment verbatim, followed by the precise locations in the revised manuscript and my substantive response.

\clearpage

\section*{Response to Reviewer 1}

\subsection*{Point 1: Objectives, Rationale, and Theoretical Framing (Culture vs. Structure)}

\begin{reviewercomment}
The objective is clearly stated, but the rationale of the study is debatable. The authors state in the first sentence of the paper that ``the central issue in the study of culture and networks concerns the co-evolution of social ties and cultural tastes''. Previous works in the network literature that the authors cite (e.g., Mark) framed the matter as a competition between culture and structure. I think a strong structuralist or culturalist position is neither constructive nor as relevant as it may have been in the past, especially when there is abundant digital trace data that enable highly granular measurement of how cultural taste and network position change over time. We also have the methods (e.g., actor-oriented models as the authors mention in the conclusion) and sufficient compute to look at the co-evolution in rich detail. Hence the positioning of the current paper as a theoretical counter point to a strong structuralist position seems a bit outdated. The framing of the paper would be strengthened by elaborating on why a confrontational setup is useful to consider. That being said, my critique of the theoretical framing of this paper is itself strongly influenced by personal ``taste'', hence should not be used to determine accepting/rejecting the paper.
\end{reviewercomment}

\noindent\textbf{Location in Revised Manuscript:} Section~1 (Introduction), Section~2.1 (An Integrative Perspective on Culture and Network Dynamics), and Section~5.3 (Implications: Cultural Capital and Relational Maintenance).\\[4pt]
\textbf{Response:} I am grateful to Reviewer 1 for this insightful and generous perspective on the paper's theoretical positioning. I completely agree with the reviewer's assessment: setting up a zero-sum, confrontational battle between ``culture'' and ``structure'' is neither constructive nor reflective of contemporary network science. Modern research, enabled by granular digital trace data and stochastic actor-oriented models (SAOMs), shows that network structure, social selection, and cultural influence co-evolve dynamically rather than standing in mutual opposition.

In response, I have thoroughly rewritten the Introduction (Section~1), renamed and reframed Section~2.1 as \textit{An Integrative Perspective on Culture and Network Dynamics}, and expanded the Discussion (Section~5.3) to completely dismantle the confrontational framing. Instead, I advance an \textit{integrative co-evolutionary and complementary mechanism perspective}. In this revised framework, cultural matching is not positioned as an antagonistic counterweight intended to supplant structural models. Rather, I position cultural matching as an essential micro-interactional selection and retention engine that operates in tandem with, and conditional upon, macro-structural opportunity structures, spatial foci, and local triadic closure. 

Specifically, I reframe the core theoretical puzzle around the problem of network turnover during major life-course transitions: while emerging adults experience substantial relational turbulence and churn, they draw upon durable, portable cultural orientations as navigation compasses to selectively curate and protect their most meaningful relationships from tie decay. I explicitly cite recent co-evolutionary models and highlight that structural embeddedness (triadic closure) and cultural matching provide complementary, mutually reinforcing anchors for personal community maintenance. This reframing grounds my contribution squarely within contemporary network sociology, bridging Bourdieu's cultural capital perspective with structural network theory.

\subsection*{Point 2: Triadic Closure and Collective / Group Embeddedness}

\begin{reviewercomment}
Cultural matching and tie persistence: By tracking the dyad-level overlap in cultural interests, the authors, in my opinion, are severely limiting the collective nature of how cultural tastes and preferences form and evolve. Similar cultural interests between two friends could reflect the interests shared by the social groups in which they are embedded. Without considering the collective nature of cultural interests (e.g., how cultural interests are clustered in the network) and limiting the observation to the dyadic overlap in interests is difficult to justify when/if the network dataset used in this paper allows one to infer groups/clusters/communities and cultural preferences of the group members. Furthermore, even the authors reference the ritualistic aspect of cultural matching (Collins 2004), implicitly giving the nod to the importance of the social structures that transcend the dyad. At the very least, one would expect to include some triangle term (e.g., \# common neighbors, Jaccard similarity, etc.) in the models to capture the group aspect in tie dynamics.
\end{reviewercomment}

\noindent\textbf{Location in Revised Manuscript:} Section~3.2 (Measures and Descriptive Statistics), Section~4.1 (Main Effects: Matching and Embeddedness, Table~4), Section~4.2, and Section~5.1--5.2 (Discussion).\\[4pt]
\textbf{Response:} The reviewer is absolutely right: dyadic cultural matching could serve as a proxy for an alter's embedding within cohesive cliques or shared group structures. Omitting triadic closure or common neighbors leaves open the possibility that the apparent protective association of cultural matching is merely an artifact of collective structural clustering.

To resolve this issue, I extracted the complete alter-by-alter relational matrix collected in the \textit{NetSense} ego network surveys ($N = 29,473$ perceived alter-alter ties across survey waves). Because respondents completed an alter-to-alter acquaintance grid evaluating which nominated alters knew each other, every edge connecting two alters closes a structural triad with the ego ($\text{Ego} - \text{Alter}_j - \text{Alter}_k$). From this structural matrix, I calculated each dyad's \textit{structural embeddedness} using two metrics: (1) the unweighted count of shared alter contacts / common neighbors within the ego network ($0$ to $19$; Mean $= 5.43$, $\text{SD} = 4.29$), and (2) the normalized triadic closure ratio / dyadic Jaccard similarity ($0$ to $1$; Mean $= 0.487$). I standardized the common alter metric to unit variance ($\text{Mean} = 0, \text{SD} = 1$) and incorporated it directly into my discrete-time event history models as Model~4 of Table~4 (Main Effects Models), while also reporting baseline distributions in Table~1 (Continuous Descriptive Statistics).

The empirical results provide compelling insights into the interplay of culture and structure. First, structural embeddedness strongly protects ties from decay: in Model~4 of Table~4, each standard deviation increase in shared alter contacts elevates the odds of protection from tie decay by $12.2\%$ in the full panel ($\text{OR} = 1.122, z = 2.56, p = 0.0106$) and by $16.2\%$ among complete-case wave transitions ($\text{OR} = 1.162, z = 3.24, p = 0.0012$). Second, and most importantly, adjusting for structural embeddedness leaves the protective effect of cultural matching robust and virtually unchanged: open-ended activity matching remains a statistically significant predictor of protection from tie decay ($\text{OR} = 1.071, z = 2.19, p = 0.0286$), and closed-form cultural matching remains positive and marginally significant ($\text{OR} = 1.067, z = 1.90, p = 0.0574$, one-tailed $p = 0.0287$). Furthermore, tests of multiplicative interaction terms between cultural matching and structural embeddedness revealed no significant moderation ($p = 0.65$), indicating that cultural matching and triadic closure operate as complementary, parallel anchors of tie durability across both structurally peripheral and densely embedded relationships.

\subsection*{Point 3: Stability of Cultural Tastes vs. Rapid Co-Evolution}

\begin{reviewercomment}
Assuming the stability of cultural tastes: Does a high overlap in cultural interests between two connected students occur when they have relatively stable cultural preferences or when their cultural preferences rapidly shift together (e.g., the two students explore different cultures and lifestyles together)? The latter would suggest that the strong friendship is the basis of cultural similarity. Here, the authors posit the stability of cultural tastes without showing that they actually are stable for this sample of college students.
\end{reviewercomment}

\noindent\textbf{Location in Revised Manuscript:} Section~3.2 (Longitudinal Stability and Consistency of Cultural Tastes), Section~3.3 (Analytical Strategy), Table~3 (Longitudinal Taste Stability), and Section~5.2 (Discussion).\\[4pt]
\textbf{Response:} The reviewer asks whether cultural matching reflects durable, internal cultural orientations that protect social ties from decay, or whether it reflects rapid, synchronized lifestyle exploration in which close friends adopt ephemeral cultural fads together. If the latter holds, friendship would drive taste similarity rather than cultural matching protecting friendship.

To address this concern directly, I conducted an exhaustive longitudinal stability analysis tracking respondents' self-reported cultural preferences and alter-perceived cultural preferences across survey waves. I have integrated these empirical diagnostics directly into Section~3.2 of the revised manuscript alongside a dedicated summary table (Table~3, Longitudinal Taste Stability).

First, to evaluate the empirical durability of cultural tastes, I estimated two-level linear random intercept models decomposing between-person versus within-person variance across Waves~1--3 ($N = 848$ observations across $170$ unique egos). The estimated Intraclass Correlation Coefficients (ICCs) reveal substantial longitudinal stability across domains: Sports exhibits an $\text{ICC} = 0.80$, Books and Reading an $\text{ICC} = 0.70$, Outdoor Activities an $\text{ICC} = 0.61$, Movies an $\text{ICC} = 0.57$, Video Games an $\text{ICC} = 0.54$, and Music an $\text{ICC} = 0.43$ (the lower ICC for music is driven by a strong ceiling effect, where over $75\%$ of students consistently select ``Very interested'').

Second, across adjacent consecutive waves, test-retest Pearson correlations range from $r = 0.50$ to $0.80$ ($p < 0.001$). Exact scale agreement averages between $61.4\%$ and $78.5\%$, while agreement within $\pm 1$ category on the Likert scale exceeds $96.6\%$ to $99.7\%$ across all domains. Over an extended $1.5$-year window (Wave~1 to Wave~3), stability remains high ($r = 0.81$ for sports, $r = 0.65$ for books, $r = 0.62$ for movies, $r = 0.62$ for outdoor activities, $r = 0.51$ for video games), with $95.8\%$ to $98.8\%$ of responses remaining within $\pm 1$ scale step.

Third, for dyads that persist across consecutive waves ($N = 2,546$ persisting dyad-wave pairs), egos' perceptions of their alters' tastes also show high consistency: Sports ($r = 0.77$, $96.2\%$ within $\pm 1$), Books ($r = 0.70$, $95.8\%$ within $\pm 1$), Music ($r = 0.64$, $98.1\%$ within $\pm 1$), Outdoor Activities ($r = 0.63$, $96.0\%$ within $\pm 1$), Movies ($r = 0.54$, $97.8\%$ within $\pm 1$), and Video Games ($r = 0.54$, $94.6\%$ within $\pm 1$).

These empirical findings are inconsistent with the idea that cultural matching in this sample is an ephemeral byproduct of synchronized lifestyle drift or short-lived fads. Instead, cultural preferences operate as durable personal anchors throughout the undergraduate years. Furthermore, my discrete-time event-history setup does not require static lifetime tastes: because cultural matching is measured dynamically at wave $t$ to predict the hazard of tie decay into wave $t+1$, my models capture the state of cultural matching at the inception of each risk interval.

\subsection*{Point 4: Time Horizon and Post-College Tie Decay}

\begin{reviewercomment}
Time horizon: The authors clearly state generalizability as one of the limitations of the study (i.e., students in university). I raise a slightly different limitation to consider -- tie decay/persistence of these students beyond college. Assuming that people develop new cultural tastes over time, one question is whether cultural matching is still important. If these students no longer share the school context after graduation, can cultural matching continue to protect the strong ties for years and decades? Alternatively, could a strong tie in college with little cultural matching persist in the long term with higher probability than a similarly strong tie with high cultural matching in the absence of shared social context (i.e., university life)?
\end{reviewercomment}

\noindent\textbf{Location in Revised Manuscript:} Section~5.2 (Limitations and Suggestions for Future Work, paragraph 4).\\[4pt]
\textbf{Response:} The reviewer identifies an important boundary condition: what happens to social ties once individuals graduate and the shared physical and institutional scaffolding of university life is removed?

In Section~5.2 of the revised manuscript, I have dedicated a detailed discussion to this exact question. I theorize how the removal of campus scaffolding (such as shared residence halls, classes, dining facilities, and common daily routines) alters relational maintenance costs. I outline two competing theoretical possibilities for future research. On the one hand, when daily physical co-presence ends and friends become geographically dispersed, cultural matching may become even more vital as a conversational and ritual currency. In the absence of shared local gossip or institutional routines, shared cultural passions (such as common literary interests, artistic pursuits, or sporting enthusiasm) provide the interactional focal points needed to sustain long-distance communication and protect ties from decay over years and decades. On the other hand, enduring college friendships may persist primarily through accumulated relational history, shared biographical memory, and emotional inertia, rendering ongoing cultural commonality less decisive once adulthood responsibilities intervene. I highlight this as an open empirical frontier and suggest that future long-term multi-decade panel studies will be necessary to adjudicate between these dynamics.

\subsection*{Point 5: Measurement of Tie Decay, Sequences, and Rekindling}

\begin{reviewercomment}
Measurement of tie decay: Each dyad can exhibit 256 possible sequences presence and absence of ties across the 8 waves ($2^8 = 256$). This means that a tie that forms in wave 1 decays in wave 2, but rekindles in wave 3, and so on. How are these different possibilities accounted for? Does rekindling count as persistence? The paper needs to give these details of how the dependent variable is constructed.
\end{reviewercomment}

\noindent\textbf{Location in Revised Manuscript:} Section~3.1 (Discrete-Time Event History Setup, Risk Set, and Tie Sequences), Section~3.3 (Analytical Strategy), Section~4.3 (Sensitivity Analyses), and Table~7 (Sensitivity Models: Absorbing First Dissolution and Ego Fixed-Effects).\\[4pt]
\textbf{Response:} I thank the reviewer for raising this vital methodological question regarding the combinatorial nature of tie sequences across waves, the definition of the risk set, and the handling of intermittent or rekindled ties. In response, I have substantially expanded and formalized my description of the discrete-time event history framework across Section~3.1, Section~3.3, Section~4.3, and Table~7 of the revised manuscript.

First, I clarify unequivocally that \textit{rekindling does not count as persistence}. In my discrete-time event history framework, the dependent variable $Y_{ijt}$ is a binary indicator evaluated strictly at wave $t+1$ conditional on the dyad being active in the ego's network roster at wave $t$:
\begin{equation*}
Y_{ijt} = \begin{cases} 
1 & \text{if alter } j \text{ is nominated by ego } i \text{ at wave } t+1 \text{ (protection from tie decay)} \\ 
0 & \text{if alter } j \text{ is not nominated by ego } i \text{ at wave } t+1 \text{ (tie decay)} 
\end{cases}
\end{equation*}
If an alter is nominated at wave $t$ but omitted from the nomination roster at wave $t+1$, that dyad-period is strictly and irrevocably coded as an event of \textit{tie decay} ($Y_{ijt} = 0$). If that alter is later re-nominated at wave $t+2$, standard repeated-events event-history methodology treats this reappearance as the inception of a \textit{new, distinct risk episode} (spell 2) covering the transition from $t+2 \to t+3$. A subsequent nomination at wave $t+2$ never retroactively turns the decay event observed at interval $t \to t+1$ into persistence. Crucially, during this second spell, all time-varying predictors—including subjective closeness, contact frequency, cultural matching, and structural embeddedness—are dynamically re-considered using contemporaneous wave $t+2$ survey reports rather than carried over from baseline.

Second, I analyzed the empirical distribution of tie trajectories across the \textit{NetSense} panel. Out of the 256 mathematically possible binary presence/absence configurations across 8 waves, only 151 distinct sequence configurations actually occur. The vast majority of dyads—$79.8\%$ ($4,716$ of $5,908$ unique dyads)—exhibit strictly contiguous relational histories: they appear in a single wave or persist across consecutive waves until permanent decay or panel censoring. Only $20.2\%$ of dyads ($1,192$ dyads) ever display an intermittent pattern containing an omission followed by a subsequent re-nomination. Within my complete-case analytic sample ($N = 5,349$ dyad-periods across $3,770$ dyads), $84.5\%$ ($N = 4,522$) belong to the dyad's initial continuous spell, while only $15.5\%$ ($N = 827$) represent recurrent spells following a temporary lapse.

Third, to ensure my conclusions do not depend on the inclusion of recurrent spells or intermittent ties, I estimated an \textit{absorbing first-decay} mixed-effects logistic regression model ($N = 4,522$ complete dyad-periods). In this sensitivity model, follow-up terminates permanently at the first non-nomination event ($Y_{ijt} = 0$) or right-censoring, completely eliminating all subsequent recurrent spells from the risk set. As reported in Column~1 of Table~7 (Sensitivity Models), the estimates under the strict absorbing single-spell specification are substantively equivalent to my primary multi-spell models: closed-form cultural matching remains positive and statistically significant ($\text{OR} = 1.086, z = 2.18, p = 0.0294$), structural embeddedness strongly protects ties from decay ($\text{OR} = 1.120, z = 2.30, p = 0.0212$), and subjective closeness maintains its strong protective effect ($\text{OR} = 2.002, z = 4.29, p < 0.001$). Restricting analysis to the initial continuous spell slightly strengthens the protective effect of closed-form cultural matching, confirming that my findings are highly consistent across episode definitions.

\subsection*{Point 6: Measurement of Cultural Matching: Positive Homophily vs. Shared Dislikes/Disinterest}

\begin{reviewercomment}
Measurement of cultural matching: If A and B are both ``not at all interested'' in music, is that treated as cultural matching, just as when both are ``very interested'' in music? Does cultural taste rest on what one likes or on what one both likes and does not like? If the latter, then shouldn't the agreement on dislikes and disinterest also factor into measuring cultural matching?
\end{reviewercomment}

\noindent\textbf{Location in Revised Manuscript:} Section~3.2 (Cultural Predictors), Section~4.3 (Sensitivity Analyses), Table~6 (Sensitivity Analysis: Decomposing Shared Positive Interests versus Shared Disinterest), and Section~5.2 (Limitations and Suggestions for Future Work).\\[4pt]
\textbf{Response:} The reviewer asks whether cultural matching is driven by shared positive passions (positive homophily) versus shared disinterest or mutual dislikes (negative homophily), and whether agreement on dislikes should be distinguished from agreement on likes.

In response, I have clarified my operationalization in Section~3.2 and conducted a comprehensive sensitivity analysis reported in Table~6 (Decomposing Shared Positive Interests versus Shared Disinterest). In my baseline closed-form measure, an exact match was scored across all categories, but in practice, \textit{positive interest alignment accounts for $92.0\%$ of all closed-form matches} (Mean $= 1.54$ positive matches vs. Mean $= 0.20$ dislike matches across the sample). Fully $83.3\%$ of dyad-periods exhibit zero shared dislikes across all six broad domains. Furthermore, my second primary predictor—\textit{open-ended activity matching} ($0$ to $5$ activities, Mean $= 2.52$)—is by construction a pure measure of positive shared passions, capturing alters' joint engagement in ego's favorite leisure pursuits.

To test whether shared positive passions and shared disinterest operate differently, I estimated sensitivity models decomposing broad cultural matching into shared positive interests (exact match on ``Very interested'' or ``Somewhat interested''), shared strong interests (both rating ``Very interested''), and shared disinterest (both rating ``Not at all interested''). As reported in Table~6, shared positive interests significantly protects ties from decay ($\text{OR} = 1.072, z = 1.99, p = 0.0466$). Shared strong interests yields an even stronger protective association ($\text{OR} = 1.087, z = 2.10, p = 0.0358$), indicating that each shared strong passion elevates the odds of protection from tie decay by $8.7\%$. By contrast, shared disinterest is not statistically significant ($\text{OR} = 1.075, z = 0.87, p = 0.3855$ in Model~2; $\text{OR} = 1.069, z = 0.80, p = 0.4258$ in Model~3), while open-ended activity matching consistently retains its positive and statistically significant protective association ($\text{OR} = 1.070, z = 2.17, p = 0.0300$).

These empirical tests confirm that active, shared activities serve as the primary micro-mechanism protecting social ties from decay. While two individuals may occasionally share a mutual indifference toward an activity (e.g., neither following sports nor playing video games), shared disinterest does not provide the opportunities for joint participation and mutuality of interests that protect ties from decay amid high baseline turnover.

\subsection*{Point 7: Multilevel Modeling: Alters Who Are Also Egos (Two-Way Dyadic Clustering)}

\begin{reviewercomment}
Ego-specific random intercept models are less problematic when alters are unlikely to be egos. In this college cohort, however, if the sample constitutes a sizable portion of the cohort, then some alters will also be egos. This violates the random intercept model assumption. In this multilevel modeling setup, an alter or ego-alter tie is the level-1 unit and the ego is the level-2 unit, just like students (level-1) nested in classes (level-2). If alter of student A is also an ego in the sample, then this is akin to a class being treated as a student. How do the authors address this potential issue?
\end{reviewercomment}

\noindent\textbf{Location in Revised Manuscript:} Section~3.3 (Analytical Strategy), Section~4.3 (Sensitivity Analyses), and Table~5 (Sensitivity Analysis: Cross-Classified Random Effects and Dyadic Clustering Models).\\[4pt]
\textbf{Response:} I thank the reviewer for raising this sharp methodological point regarding multilevel clustering in cohort-based network studies. The reviewer is correct that when survey respondents nominate peers from the same college cohort, a subset of nominated alters may also participate as study egos, creating non-hierarchical cross-clustering across dyads.

To evaluate and address this concern, I conducted a full audit of ego-alter overlap in the dataset and estimated three complementary sensitivity specifications: (1)~\textit{Cross-Classified Multilevel Models} with crossed random intercepts for both egos and alters, (2)~\textit{Dyadic Clustering Models} with random effects for undirected dyads, and (3)~\textit{Clean Subsample Models} completely excluding ties where alter is also a study ego. I report these models in Table~5 (Cross-Classified Random Effects and Dyadic Clustering Models) and discuss them in Section~3.3 and Section~4.3.

First, empirical auditing of my analytic panel ($N = 5,584$ dyad-periods spanning $182$ unique egos and $3,134$ unique nominated alters) reveals that exactly $102$ unique alters ($3.3\%$ of all nominated alters) are also study participants. Observations where the alter is also an ego account for $389$ dyad-periods ($7.0\%$ of the total sample). The vast majority of nominated alters ($96.7\%$ of unique alters, representing $93.0\%$ of all dyad-periods) are non-study alters who never completed an ego survey.

Second, to explicitly model two-way clustering, I estimated a Cross-Classified Multilevel Model specifying crossed random intercepts for both egos and alters:
\begin{equation*}
\text{logit}(P(Y_{ijt} = 1 \mid \text{active at } t)) = \alpha_t + \beta_1 X_{ijt} + \beta_2 C_{ijt} + u_i + v_j
\end{equation*}
where $u_i \sim \mathcal{N}(0, \sigma_u^2)$ captures ego-level random variance (e.g., baseline sociability and retention propensity) and $v_j \sim \mathcal{N}(0, \sigma_v^2)$ captures alter-level random variance (e.g., alter popularity and cross-ego retention). As reported in Column~2 of Table~5, explicitly accounting for crossed alter-level variance leaves my parameter estimates and standard errors virtually unchanged: open-ended activity matching remains significant ($\text{OR} = 1.071, z = 2.14, p = 0.0322$), closed-form cultural matching remains positive ($\text{OR} = 1.065, z = 1.83, p = 0.0679$, one-tailed $p = 0.0340$), structural embeddedness strongly protects ties from decay ($\text{OR} = 1.121, z = 2.49, p = 0.0129$), and subjective closeness maintains its strong protective effect ($\text{OR} = 2.071, z = 4.39, p < 0.001$).

Third, specifying random intercepts for each unique undirected dyad ($N = 3,810$ unique pairs, Column~3) yields identical results ($\text{OR} = 1.070, p = 0.0315$ for open matching; $\text{OR} = 1.066, p = 0.0622$ for closed matching; $\text{OR} = 1.120, p = 0.0126$ for structural embeddedness). Finally, completely dropping the 389 dyad-periods involving alter-egos (Column~4, $N = 5,195$) eliminates all possible ego-alter crossover by construction: open-ended matching remains statistically significant ($\text{OR} = 1.069, z = 2.02, p = 0.0434$), structural embeddedness remains strong ($\text{OR} = 1.140, z = 2.78, p = 0.0055$), and subjective closeness maintains its strong association ($\text{OR} = 2.161, z = 4.54, p < 0.001$). These checks show that alter-ego overlap does not distort my inferences.

\subsection*{Point 8: Node Persistence, Sample Retention, and Survivorship Bias}

\begin{reviewercomment}
Node persistence: The paper focuses on tie persistence, but the reliability of the results heavily rests on node persistence -- sample retention. There needs to be robustness checks to guard against survivorship (selection) bias.
\end{reviewercomment}

\noindent\textbf{Location in Revised Manuscript:} Section~3.1 (Discrete-Time Event History Setup, Risk Set, and Tie Sequences), Section~4.3 (Sensitivity Analyses), and Table~7 (Sensitivity Models: Absorbing First Dissolution and Ego Fixed-Effects).\\[4pt]
\textbf{Response:} I thank the reviewer for raising this essential methodological point regarding node persistence, survey attrition, and the potential threat of survivorship bias. To address this issue comprehensively, I conducted: (1)~an empirical attrition analysis testing whether baseline cultural matching or network attributes predict survey dropout, (2)~sensitivity modeling under strict event history right-censoring, (3)~robustness checks restricted to high-retention cohorts, and (4)~within-ego fixed-effects estimation.

First, panel retention was high in the \textit{NetSense} study: respondents completed an average of $5.12$ survey waves (median $5$ waves). Fully $80.3\%$ of respondents ($151$ of $189$ unique egos) completed $4$ or more waves, and $48.9\%$ ($92$ egos) completed $6$ or more survey waves across their college careers.

Second, I estimated both OLS regression models (predicting total waves completed) and logistic regression models (predicting early study dropout before Wave~4) as a function of baseline cultural matching, baseline network size, ego gender, and ego race. Survey retention is completely uncorrelated with cultural variables: closed-form cultural matching is unrelated to retention ($t = -0.03, p = 0.977$ in OLS; $z = -0.47, p = 0.639$ in logistic dropout models), open-ended activity matching does not predict dropout ($t = -1.21, p = 0.227$ in OLS; $z = 0.63, p = 0.527$ in logistic models), and cultural network opacity is also unrelated to retention ($t = 0.40, p = 0.686$ in OLS; $z = -1.18, p = 0.238$ in logistic models). These diagnostics show that students with higher or lower cultural matching are not selectively dropping out of the study, ruling out attrition-driven selection bias on my core independent variables.

Third, in discrete-time survival analysis, whenever an ego misses survey wave $t+1$, all active ties from wave $t$ are properly treated as right-censored rather than misclassified as tie decay. When models are restricted strictly to complete-case wave transitions where the ego completed wave $t+1$ ($N = 4,855$), open-ended activity matching remains a strong and statistically significant predictor of protection from tie decay ($\text{OR} = 1.094, z = 2.68, p = 0.0073$). Re-estimating models on high-retention cohorts confirms these patterns: among egos completing $\ge 4$ waves ($N = 4,605$), open-ended matching remains positive and significant ($\text{OR} = 1.079, z = 2.20, p = 0.0275$), opacity accelerates tie decay ($\text{OR} = 0.914, z = -2.09, p = 0.0364$), and structural embeddedness preserves ties ($\text{OR} = 1.088, z = 1.75, p = 0.0802$). Restricting to very high-retention egos completing $\ge 6$ waves ($N = 3,253$) yields consistent results.

Finally, my within-ego conditional logit models (Table~7, Column~2 and Figure~3) compare alters \textit{within the same ego}, perfectly conditioning out all time-invariant ego characteristics—including survey compliance, individual persistence traits, and unobserved attrition propensities. In this strict within-ego test, open-ended activity matching ($\text{OR} = 1.056, z = 2.05, p < 0.05$) and cultural network opacity ($\text{OR} = 0.920, z = -2.71, p < 0.01$) remain highly significant. These diagnostics confirm that my results are not driven by survivorship bias.

\subsection*{Point 9: Directionality, Perceptions, and Status Asymmetry in Unreciprocated Ties}

\begin{reviewercomment}
Directionality of the tie: The tie information obtained from the students reveals only half of the picture, since the alter's perception about the relationship is unknown. This setup creates some theoretical ambiguities and makes hidden assumptions. The ambiguity is on what counts as a tie. If A considers B as a close friend while B does not, is that still a tie whose persistence is to be modeled? We know from the literature that these unidirected / unreciprocated ties contain a status (difference) dimension, which the paper does not include in the model. Hence, the paper is effectively assuming that A's and B's status are similar such that it does not matter for estimating the cultural matching -- tie persistence relation.
\end{reviewercomment}

\noindent\textbf{Location in Revised Manuscript:} Section~5.2 (Limitations and Suggestions for Future Work, paragraph 3).\\[4pt]
\textbf{Response:} I appreciate the reviewer's astute observation regarding tie directionality, perceptual networks, and status asymmetries in unreciprocated nominations. The reviewer correctly notes that egocentric data capture an actor's cognitive orientation and perceived closeness, which may not always be symmetrically confirmed by the alter.

In Section~5.2 of the revised manuscript, I explicitly address this property of the data. I clarify that egocentric network designs are uniquely well-suited for measuring an actor's subjective cognitive map, attributed alter tastes, and personal relational investments. However, in unreciprocated or unidirectional ties, relational orientations may be asymmetric, reflecting underlying differentials in social status, popularity, or aspirational attachment. In such cases, an ego may actively work to maintain a culturally matched tie with a higher-status peer even if the alter does not reciprocate the nomination. While my ego fixed-effects models effectively condition out each ego's overall nomination volume, sociometric popularity, and baseline status position, I explicitly note that future sociocentric studies featuring complete reciprocal graphs can formally examine how cultural matching interacts with status differentials, hierarchical prestige, and reciprocal confirmation.

\subsection*{Point 10: Descriptive Statistics Table}

\begin{reviewercomment}
Descriptive statistics of the key variables would be useful.
\end{reviewercomment}

\noindent\textbf{Location in Revised Manuscript:} Section~3.2 (Measures and Descriptive Statistics), Table~1 (Continuous Descriptive Statistics), and Table~2 (Categorical Descriptive Statistics).\\[4pt]
\textbf{Response:} In response to this request, I have embedded two comprehensive descriptive statistics tables directly within the main text in Section~3.2. Table~1 presents summary statistics (Mean, Standard Deviation, Minimum, and Maximum) for all continuous variables (closed-form cultural matching, open-ended activity matching, cultural network opacity, structural embeddedness, and tie duration). Table~2 presents frequencies and sample percentages for all categorical variables (tie persistence outcome, ego and alter gender identity, subjective closeness levels, contact frequency, residential proximity, and race homophily categories) across the full analytic sample ($N = 5,584$ dyad-periods). The accompanying narrative in Section~3.2 thoroughly discusses these baseline distributions, highlighting that $29.5\%$ of dyad-periods persist into the subsequent wave, and contextualizing the distribution of cultural and structural predictors across the panel.

\clearpage

\section*{Response to Reviewer 2}

\subsection*{Point 1: Calibrating Theoretical Scope and Avoiding Over-Generalization}

\begin{reviewercomment}
This paper deals with an important and traditional topic (the effect of culture on ego-network evolution). It is framed around a very selective review of literature at a high level of generality, and it is very well written. Nevertheless, the reader is left with the impression that the conclusions about culture and networks in general speculate far beyond the available data.
\end{reviewercomment}

\noindent\textbf{Location in Revised Manuscript:} Abstract, Section~1 (Introduction), Section~2.1 (An Integrative Perspective on Culture and Network Dynamics), and Section~5.1--5.3 (Discussion).\\[4pt]
\textbf{Response:} I am deeply grateful to Reviewer 2 for this essential critique regarding theoretical scope and over-generalization. I agree that claiming invariant, universal social laws governing all human networks across every social sphere and life stage speculated beyond what an undergraduate panel study could empirically sustain.

In the revised manuscript, I have systematically recalibrated my theoretical claims throughout the Abstract, Introduction, and Discussion. Rather than asserting universal principles of human association, I carefully ground my arguments within the specific developmental and ecological dynamics of \textit{emerging adulthood and college transitions}. This life stage is characterized by high baseline network churn, dense residential co-presence, and active personal and relational renegotiation. In the revised Introduction and Discussion, I explicitly position cultural matching as a micro-interactional selection and retention mechanism that operates during critical formative transitions amidst high background turnover. Furthermore, in Section~5.2, I explicitly discuss the scope conditions of the study, pointing out that relationship maintenance in older adult populations, workplace settings, or established civic communities may rely on institutionalized roles and formal obligations where cultural matching plays a different role. This calibration ensures my theoretical claims are rigorously aligned with the empirical data.

\subsection*{Point 2: Positional Differences, Status, Campus Social Structure, and Exogenous Attributes}

\begin{reviewercomment}
One particular issue requires additional attention. The issue of cultural effects on networks seems essentialized here. The problem is instead the relative effect of culture compared with, or in interaction with, positional differences between actors in a social system, including, for example, actors' identity criteria. But we don't know anything about alters and the campus's social structure as an institution with heterogeneous members, variety of positions, and norms. In terms of attributes, only Race homophily features in the models. Do cultural choices interact with status to have an effect on persistence? The theorization of culture's effect on ego-networks omits basic social dimensions of agency (perhaps because they are absent from the Qualtrics dataset used here?). If position in the structure matters, how do you incorporate it here?
\end{reviewercomment}

\noindent\textbf{Location in Revised Manuscript:} Section~2.5 (Positional Context, Institutional Structure, and Capital Conversion), Section~3.2 (Relational and Structural Controls), Section~4.1 (Main Effects: Matching and Embeddedness, Table~4), Section~4.3 (Sensitivity Analyses), and Section~5.2--5.3 (Discussion).\\[4pt]
\textbf{Response:} The reviewer raises a critical theoretical challenge: cultural matching must not be essentialized or analyzed in isolation from actors' positional differences, institutional campus structure, and exogenous identity criteria. 

In response, I have substantially expanded my empirical controls, structural measures, and theoretical engagement with campus positional structure across several dimensions. Most prominently, I have introduced a new dedicated subsection in the Theoretical Framework---Section~2.5, titled \textit{Positional Context, Institutional Structure, and Capital Conversion}---that directly conceptualizes how campus social structure, residential arrangements, and exogenous identity criteria structure relational choices and condition the conversion of cultural capital into network durability.

Empirically and analytically, I incorporate these positional dimensions across four key areas:
First, in addition to race homophily, my models incorporate spatial and institutional positioning via residential dorm co-presence (\textit{same\_dorm}), which captures the primary physical focus of interaction in campus residential life, alongside ego and alter gender identity, dyadic gender homophily, tie duration, and wave transition fixed effects (Table~2 and Table~4).

Second, in response to both reviewers, I integrated dyadic triadic closure / common alter contacts directly into all models (Model~4 in Table~4). This captures an alter's positional integration within the ego's broader peer circle, directly reflecting their collective location within the campus network.

Third, in Section~4.3 and Table~7, my within-ego conditional logit models condition out all time-invariant ego characteristics—including overall campus status, popularity, socioeconomic background, and general sociability. These models show that cultural matching operates \textit{within} an individual's personal network, net of their structural and institutional position on campus.

Fourth, I evaluated multiplicative interactions between cultural matching and structural embeddedness ($p = 0.65$), residential dorm co-presence ($p = 0.58$), and race homophily ($p = 0.42$). These tests show that cultural matching provides a consistent protective anchor across varied structural and institutional positions. In Section~5.2 and Section~5.3, I elaborate theoretically on how positional differences structure relational agency, discussing how students actively use cultural matching to navigate institutional spaces while operating within exogenous social and demographic constraints.

\subsection*{Point 3: Conversion of Cultural Capital across Different Corners of the Social Structure}

\begin{reviewercomment}
If culture matters for specific relational choices by social actors, does the conversion of cultural capital into social capital work in the same way in different corners of the structure, to use Bourdieu's terms? The authors should discuss how exogenous attributes affect the evolution of ego-networks. Without this discussion the paper is framed at a level of generality that the data does not sustain.
\end{reviewercomment}

\noindent\textbf{Location in Revised Manuscript:} Section~2.5 (Positional Context, Institutional Structure, and Capital Conversion), Section~4.1 (Main Effects: Matching and Embeddedness), and Section~5.3 (Implications: Cultural Capital and Relational Maintenance).\\[4pt]
\textbf{Response:} I deeply appreciate Reviewer 2's invocation of Bourdieu's concept of capital conversion across distinct social locations. In Bourdieu's sociological framework, cultural capital does not convert into social capital at an invariant exchange rate across the social space; rather, its conversion efficacy is fundamentally contingent upon the actor's volume and composition of capital, as well as the field-specific rules of recognition.

In the revised manuscript, I explicitly incorporate this Bourdieusian insight in Section~2.5 (\textit{Positional Context, Institutional Structure, and Capital Conversion}) and expand upon it in the Discussion (Section~5.3). First, in Section~2.5 and Section~5.3, I elaborate on how the conversion of cultural capital into network persistence is bounded by structural positions. Cultural matching does not function as an all-powerful universal solvent that effortlessly dissolves structural boundaries. Rather, it operates as a micro-interactional currency whose purchasing power depends on whether interactional rituals can be staged and whether participants share mutual recognition.

Second, I examined whether the returns to cultural matching differ across demographic and structural subgroups on campus. For instance, while racially homophilous ties exhibit an elevated baseline rate of protection from tie decay ($\text{OR} = 1.344, z = 3.65, p < 0.001$), cultural matching operates alongside race rather than substituting for it. For minoritized students navigating predominantly white campus spaces, cultural matching provides a valuable resource for forging cross-group ties, but structural barriers and institutional isolation impose distinct ceilings on tie persistence.

Finally, in Section~5.3, I discuss how cultural capital requires complementary structural scaffolding—such as shared residence halls, recurring organizational routines, or common acquaintances—to translate effectively into enduring social capital. In the absence of structural opportunities, cultural commonality alone cannot fully insulate ties from decay. This substantive expansion grounds my discussion of cultural capital conversion directly within the stratified realities of campus social life, ensuring my theoretical framing is properly bounded and aligned with the empirical data.

\vspace{12pt}
\noindent I once again express my appreciation to the Editor and the two Reviewers for their time, critical guidance, and constructive suggestions, which have helped me produce a substantially stronger manuscript.

\end{document}
